\documentclass[10pt,a4paper]{amsart}
\usepackage[margin=19mm]{geometry}
\usepackage{amsmath,amssymb,mathtools}
\usepackage[scr=rsfs,cal=euler]{mathalfa}
\usepackage{xfrac}
\usepackage[unicode,hidelinks,bookmarks=false]{hyperref}
\newcommand{\ie}{i.e.\ }
\newcommand{\cf}{cf.\ }
\newcommand{\resp}{respectively\ }
\newcommand{\Subsection}{\subsection}
\providecommand{\nobreakdash}{\nobreak\hbox{-}}
\setlength{\emergencystretch}{2em}
\multlinegap0pt
\usepackage{mathtools}
\usepackage{amssymb}
\usepackage[shortlabels]{enumitem}
\usepackage{tikz}
\newtheorem{lemma}{Lemma}
\newtheorem{claim}{Claim}
\newcommand{\F}{\mathbb{F}}
\DeclareMathOperator{\diag}{diag}
\renewcommand*{\leq}{\leqslant}
\DeclareMathOperator{\nega}{neg}
\DeclareMathOperator{\supp}{supp}
\newcommand{\trim}{/\!\!/}
\title{Totally equimodular matrices: decomposition and triangulation}
\author{Patrick Chervet, Roland Grappe, Mathieu Vall\'ee}
\date{Source: arXiv:2504.05930v2 (CC BY 4.0)}
\begin{document}
\maketitle

\noindent\emph{Source and licence.} Totally equimodular matrices: decomposition and triangulation. Version arXiv:2504.05930v2 (CC BY 4.0), distributed by arXiv under the Creative Commons Attribution 4.0 International licence, http://creativecommons.org/licenses/by/4.0/, as recorded in the arXiv licence metadata for this exact version and shown on the abstract page https://arxiv.org/abs/2504.05930v2. Full licence notice: \texttt{licenses/arxiv-2504.05930-CC-BY-4.0.txt}.\par\smallskip

\noindent\textit{Editorial scope.} Counted unit: the article's lemma lemma:trimming\_and\_complement\_orbit (source0.tex lines 1833--1841). The article's complete original proof of it is delivered with it (source0.tex lines 1842--2098, one \texttt{proof} environment of 257 lines). No mathematics is written, repaired or re-derived: the statement and the proof are the article's own lines, in the article's own notation, and no retained line is substituted or re-flowed.  Retained uncounted as the source-local context the proof needs: lines 1796--1811; lines 1817--1822.  Nothing was omitted from any retained span, so the package carries no omission record.  No retained line contains a citation, so the wrapper carries no bibliography and no bibliography entry.  No retained line contains a figure environment, an image-inclusion command or drawing code, so no image is delivered and the package declares no asset dependency.  Editorial formatting only, in addition: an AMS \texttt{amsart} preamble that restates the article's own declarations and macros used by the retained lines, namely the environments \texttt{lemma}, \texttt{claim} on separate counters, together with the standard packages those lines need.  The retained environments print in this document as Lemma 1, Claim 1 in order of appearance, whereas the article prints the same items with the numbering of its own full document.  The complete native source of the article as distributed in the arXiv e-print for this exact version is bundled unchanged as \texttt{sources/176177/source0.tex}.
\begin{lemma}\label{lemma:correspondence_sign}
Let $A$ be a $\pm1$ matrix of size $m\times n$. Then, for every $\varepsilon \in \{\pm1\}^m$, we have:

$$\nega{(A\diag(\varepsilon))} =\begin{bmatrix}
\nega{(A)}_1 \oplus \nega{(\varepsilon)}^\top\\
\vdots \\
\nega{(A)}_m \oplus \nega{(\varepsilon)}^\top
\end{bmatrix},$$

and for every $\mu \in \{\pm1\}^n$, we have:

$$\nega{(\diag(\mu)A)} =\begin{bmatrix}
\nega{(A)}^1 \oplus \nega{(\mu)} &\cdots &
\nega{(A)}^n \oplus \nega{(\mu)}
\end{bmatrix}.$$
\end{lemma}

\begin{equation}\label{equation:core}
A=\begin{bmatrix}
1 & \mathbf{1}^\top \\
\mathbf{1} & \mathbf{J} - 2B
\end{bmatrix},
\end{equation}

\begin{lemma}\label{lemma:trimming_and_complement_orbit}
Let $A$ be a $\pm1$ matrix of size $(m+1)\times(n+1)$, and $B$ a core of $A$ as in~\eqref{equation:core}.
Then, for any pivot $p=(i+1,j+1)$ of $A$, with $0\leq i\leq m$ and $0\leq j \leq n$, we have
\begin{equation}\label{equation:trim_complement}
A\trim p \simeq -2A_{i+1}^{j+1}\diag(\varepsilon) B^{[j]}_{[i]}\diag(\mu),
\end{equation}
where $\diag(\varepsilon)$ and $\diag(\mu)$ are signing matrices with $\varepsilon = \mathbf{1} -2 B^j$ and $\mu=\mathbf{1}^\top -2 B_i$, and the convention $B^0=\mathbf{0}$ and $B_0=\mathbf{0}^\top$.
%	for $\diag(\varepsilon_j)$ and $\diag(\varepsilon_i)$ diagonal matrices with $\varepsilon_j=-(2B^j-\mathbf{1})$, and $\varepsilon_i=-(2B_i-\mathbf{1}^\top)$.
\end{lemma}

\begin{proof}
First, note that $A\trim(1,1) = -2B = -2A_1^1 \diag(\mathbf{1}-2B^0) B_{[0]}^{[0]}\diag(\mathbf{1}^\top-2B_0)$.
Similarly, with a matrix of the form
\begin{equation}\label{equation:nice_form}
A'=\begin{array}{c@{}c}
	& \begin{array}{ccc}
		& j & 
	\end{array}\\
	\begin{array}{c}
		\\i \\  \\
	\end{array} &\begin{bmatrix}
		\mathbf{J}-2D & \mathbf{1} & \mathbf{J}-2E\\
		\mathbf{1}^\top & 1 & \mathbf{1}^\top\\
		\mathbf{J}-2F & \mathbf{1} & \mathbf{J}-2G
	\end{bmatrix}\\&\end{array}\!\!, \text{ we have } \nega{(A')}=\begin{array}{c@{}c}
	& \begin{array}{ccc}
		& j & 
	\end{array}\\
	\begin{array}{c}
		\\i \\  \\
	\end{array} &\begin{bmatrix}
		D & \mathbf{0} & E\\
		\mathbf{0}^\top & 0 & \mathbf{0}^\top\\
		F & \mathbf{0} & G
	\end{bmatrix}\\&\end{array}\!\!,
	\end{equation}
	and we get $A'\trim(i,j) = -2\begin{bmatrix}
D&E\\F&G
\end{bmatrix}$.
\begin{claim}
For $A\in\{\pm1\}^{m\times n}$, $\varepsilon\in\{\pm1\}^m$, $\mu\in\{\pm1\}^n$, $i\in\{1,\ldots,m\}$, and $j\in\{1,\ldots,n\}$, we have:
\begin{equation}\label{equation:diag_trim_col}
	\left(\diag(\varepsilon)A\diag(\mu)\right)\trim(i,j) = \diag(\varepsilon_{\widehat{i}})(A\trim(i,j))\diag(\mu_{\widehat{j}}).
\end{equation}

\end{claim}
\begin{proof}
Since $A$ is $\pm1$, we have $A_i^j =\frac{1}{A_i^j}$, and the classical formula for pivoting yields 
$A\trim (i,j) = A_{\widehat{i}}^{\widehat{j}} - A_i^j A_{\widehat{i}}^{j} A_{i}^{\widehat{j}}$.
By applying this formula to $M=\diag(\varepsilon)A\diag(\mu)$, and since $\varepsilon$ and $\mu$ are $\pm1$, we have:
\begin{align*}
	M\trim (i,j) &= M_{\widehat{i}}^{\widehat{j}} - M_i^j M_{\widehat{i}}^{j} M_{i}^{\widehat{j}} \\ &= \diag(\varepsilon_{\widehat{i}})(A_{\widehat{i}}^{\widehat{j}})\diag(\mu_{\widehat{j}}) - \varepsilon_iA_i^j\mu_j (\diag(\varepsilon_{\widehat{i}})A_{\widehat{i}}^{j} \mu_j)(\varepsilon_i A_{i}^{\widehat{j}}\diag(\mu_{\widehat{j}}))\\
	&=\diag(\varepsilon_{\widehat{i}})\left(A_{\widehat{i}}^{\widehat{j}} - A_i^j A_{\widehat{i}}^{j} A_{i}^{\widehat{j}}\right)\diag(\mu_{\widehat{j}})\\
	&= \diag(\varepsilon_{\widehat{i}})(A\trim(i,j))\diag(\mu_{\widehat{j}}).
\end{align*}
Thus, the claim is proved.
\end{proof}
Our strategy is to sign $A$ in order to obtain a matrix $A'$ of the form~\eqref{equation:nice_form}, and then perform the $(i,j)$-trim, which is the matrix $\nega{(A')}$ after removing the $i$-th row and $j$-th column of $\mathbf{0}$, and multiplying it by $-2$.
We then compare it to the different types of matrices in the complement orbit of $B$.
\paragraph*{\textbf{Case 1: $(i+1,1)$-pivot and $(1,j+1)$-pivot}}~\\
Let  $1\leq i\leq m$, by Lemma~\ref{lemma:correspondence_sign}, we have
$$\nega{(A\diag(A_{i+1}))} =\begin{bmatrix}
\nega{(A)}_1 \oplus \nega{(A)}_{i+1}\\
\nega{(A)}_2 \oplus \nega{(A)}_{i+1}\\
\vdots \\
\nega{(A)}_{i} \oplus \nega{(A)}_{i+1}\\
\nega{(A)}_{i+1} \oplus \nega{(A)}_{i+1}\\
\nega{(A)}_{i+2} \oplus \nega{(A)}_{i+1}\\
\vdots \\
\nega{(A)}_m \oplus \nega{(A)}_{i+1}
\end{bmatrix} = \begin{bmatrix}
0&\mathbf{0}^\top \oplus B_i\\
0& B_1 \oplus B_i\\
\vdots & \vdots \\
0& B_{i-1}\oplus B_i\\
0& \mathbf{0}^\top\\
0& B_{i+1} \oplus B_i\\
\vdots & \vdots \\
0 & B_m \oplus B_i
\end{bmatrix},$$
Since $\mathbf{0}^\top \oplus B_i=B_i$, we have $(A\diag(A_{i+1}))\trim (i+1,1) \simeq -2B_{[i]}$, up to a cyclic permutation of the first $i$ rows. By~\eqref{equation:diag_trim_col} and $A_{i+1} = 1^\top -2B_i$ we obtain: $$A\trim(i+1,1)\simeq-2B_{[i]}\diag(\mathbf{1}^\top-2B_i)\simeq-2 A_{i+1}^1 B_{[i]}\diag(\mathbf{1}^\top-2B_i).$$
We obtain by similar computations and by~\eqref{equation:diag_trim_col} that, for $1\leq j\leq n$: $$A\trim(1,j+1)\simeq-2\diag(\mathbf{1}-2B^j)B^{[j]}\simeq-2A_1^{j+1}\diag(\mathbf{1}-2B^j)B^{[j]}.$$
\medskip

\paragraph*{\textbf{Case 2: $(i+1,j+1)$-pivots}}~\\
Now, let $1\leq i\leq m$ and $1\leq j\leq n$. We first look at how $B_{[i]}^{[j]}$ looks like, depending on $B_i^j$.
\renewcommand{\arraystretch}{1.5}
Let $J_1=\supp(B_i)\setminus \{j\}$, $J_0 =\{1,\ldots,m\}\setminus \supp(B_i)$, $I_1=\supp(B^j)\setminus\{i\}$ and $I_0 =\{1,\ldots,n\}\setminus \supp(B^j)$.
Up to putting row $i$ at the top of $B$ and column $j$ to the left, and up to reorganizing the rows and columns, and if $B_i^j=0$, we may assume that
$$B=\begin{array}{c@{}c}
& \begin{array}{lcc}
	j\;\; & J_1\;\;& J_0
\end{array}\\
\begin{array}{c}
	i\\I_1 \\ I_0
\end{array} &\begin{bmatrix}
	0 &\mathbf{1}^\top & \mathbf{0}^\top\\
	\mathbf{1} & B_{I_1}^{J_1} & B_{I_1}^{J_0} \\
	\mathbf{0} & B_{I_0}^{J_1} & B_{I_0}^{J_0} \\
\end{bmatrix}\\&\end{array}\!\!, \text{ therefore } B_{[i]}^{[j]}=\begin{array}{c@{}c}
& \begin{array}{lcc}
	j\;\; & J_1\;\; & J_0
\end{array}\\
\begin{array}{c}
	i\\I_1 \\ I_0
\end{array} &\begin{bmatrix}
	0 &\mathbf{1}^\top & \mathbf{0}^\top\\
	\mathbf{1} & B_{I_1}^{J_1} & \overline{B_{I_1}^{J_0}} \\
	\mathbf{0} & \overline{B_{I_0}^{J_1}} & B_{I_0}^{J_0} \\
\end{bmatrix}\\&\end{array}.$$
If $B_i^j=1$, we have 
$$B=\begin{array}{c@{}c}
& \begin{array}{lcc}
	j\;\; & J_1\;\; & J_0
\end{array}\\
\begin{array}{c}
	i\\I_1 \\ I_0
\end{array} &\begin{bmatrix}
	1 &\mathbf{1}^\top & \mathbf{0}^\top\\
	\mathbf{1} & B_{I_1}^{J_1} & B_{I_1}^{J_0} \\
	\mathbf{0} & B_{I_0}^{J_1} & B_{I_0}^{J_0} \\
\end{bmatrix}\\&\end{array}\!\!, \text{ therefore } B_{[i]}^{[j]}=\begin{array}{c@{}c}
& \begin{array}{ccc}
	j\;\; & J_1\;\; & J_0
\end{array}\\
\begin{array}{c}
	i\\I_1 \\ I_0
\end{array} &\begin{bmatrix}
	1 &\mathbf{0}^\top & \mathbf{1}^\top\\
	\mathbf{0} & \overline{B_{I_1}^{J_1}} & B_{I_1}^{J_0} \\
	\mathbf{1} & B_{I_0}^{J_1} & \overline{B_{I_0}^{J_0}} \\
\end{bmatrix}\\&\end{array}.$$

We now relate the two cases $B_i^j=0$ and $B_i^j=1$ with the cases $A_{i+1}^{j+1}=1$ and $A_{i+1}^{j+1}=-1$, respectively.

Assume $A_{i+1}^{j+1}=1$, and let $J_{-} = \supp(\nega(A_{i+1}))$ and $J_{+} = (\{2,\ldots,m+1\}\setminus \supp(\nega(A_{i+1})))\setminus\{j+1\}$ and 
$I_{-} = \supp(\nega(A^{j+1}))$ and $I_{+} = (\{2,\ldots,n+1\}\setminus \supp(\nega(A^{j+1})))\setminus\{i+1\}$.
We have, up to reordering the rows and columns: 
$$A \simeq \begin{array}{c@{}c}
& \begin{array}{cccc}
	1 & J_-& j+1 & J_+
\end{array}\\
\begin{array}{c}
	1\\I_- \\ i+1 \\ I_+
\end{array} &\begin{bmatrix}
	1& \mathbf{1}^\top &1&\mathbf{1}^\top\\ 
	\mathbf{1} & A_{I_-}^{J_-} & \mathbf{-1} &  A_{I_-}^{J_+}\\
	1 & \mathbf{-1}^\top & 1 & \mathbf{1}^\top \\
	\mathbf{1}& A_{I_+}^{J_-} & \mathbf{1} & A_{I_+}^{J_+}
\end{bmatrix}\\&\end{array},$$
thus,
$$\diag(A^{j+1})A\diag(A_{i+1}) \simeq \begin{array}{c@{}c}
& \begin{array}{cccc}
	1 & J_-& j+1 & J_+
\end{array}\\
\begin{array}{c}
	1\\I_- \\ i+1 \\ I_+
\end{array} &\begin{bmatrix}
	1 & \mathbf{-1}^\top & 1 & \mathbf{1}^\top \\
	\mathbf{-1} & A_{I_-}^{J_-} & \mathbf{1} & - A_{I_-}^{J_+}\\
	1& \mathbf{1}^\top &1&\mathbf{1}^\top\\ 
	\mathbf{1}& -A_{I_+}^{J_-} & \mathbf{1} & A_{I_+}^{J_+}
\end{bmatrix}\\&\end{array},$$
therefore,
$$\nega (\diag(A^{j+1})A\diag(A_{i+1})) \simeq\begin{bmatrix}
0 &\mathbf{1}^\top & 0& \mathbf{0}^\top\\
\mathbf{1} & B_{I_1}^{J_1} &\mathbf{0}& \overline{B_{I_1}^{J_0}} \\
0 & \mathbf{0}^\top & 0 &\mathbf{0}\\
\mathbf{0} & \overline{B_{I_0}^{J_1}} &\mathbf{0}& B_{I_0}^{J_0} \\
\end{bmatrix},$$
with $J_1 = \{j-1\colon j\in J_-\}$, $J_0 = \{j-1\colon j\in J_+\}$, $I_1 = \{i-1\colon i\in I_-\}$, $I_0 = \{i-1\colon i\in I_+\}$.
This yields, $\diag(A^{j+1})A\diag(A_{i+1})\trim(i+1,j+1) \simeq -2 B_{[i]}^{[j]}$, up to cyclic permutation of the first $i$ rows and $j$ columns.
Finally, by~\eqref{equation:diag_trim_col}
\begin{align*}
A\trim(i+1,j+1) &\simeq -2 \diag(\mathbf{1} -2 B^j) B_{[i]}^{[j]}\diag(\mathbf{1}-2 B_i)\\
& \simeq -2 A_{i+1}^{j+1} \diag(\mathbf{1} -2 B^j) B_{[i]}^{[j]}\diag(\mathbf{1}-2 B_i).
\end{align*}

Assume $A_{i+1}^{j+1}=-1$, and let $J_{-} = \supp(\nega(A_{i+1}))\setminus\{j+1\}$ and $J_{+} = \{2,\ldots,m+1\}\setminus \supp(\nega(A_{i+1}))$ and 
$I_{-} = \supp(\nega(A^{j+1}))\setminus\{i+1\}$ and $I_{+} = \{2,\ldots,n+1\}\setminus \supp(\nega(A^{j+1}))$, we have:
$$A = \begin{array}{c@{}c}
& \begin{array}{cccc}
	1 & J_-& j+1 & J_+
\end{array}\\
\begin{array}{c}
	1\\I_- \\ i+1 \\ I_+
\end{array} &\begin{bmatrix}
	1 & \mathbf{1}^\top & 1 & \mathbf{1}^\top \\
	\mathbf{1} & A_{I_-}^{J_-} & \mathbf{-1} & A_{I_-}^{J_+}\\
	1& \mathbf{-1}^\top &-1&\mathbf{1}^\top\\ 
	\mathbf{1}& A_{I_+}^{J_-} & \mathbf{1} & A_{I_+}^{J_+}
\end{bmatrix}\\&\end{array}.$$
%then,
%$$A\diag(A_{i+1}) = \begin{array}{c@{}c}
%	& \begin{array}{cccc}
	%	1 & J_-& j+1 & J_+
	%	\end{array}\\
%	\begin{array}{c}
	%	1\\I_- \\ i+1 \\ I_+
	%	\end{array} &\begin{bmatrix}
	%		1 & \mathbf{-1}^\top & -1 & \mathbf{1}^\top \\
	%		\mathbf{1} & -A_{I_-}^{J_-} & \mathbf{1} & A_{I_-}^{J_+}\\
	%		1& \mathbf{1}^\top &1&\mathbf{1}^\top\\ 
	%		\mathbf{1}& -A_{I_+}^{J_-} & \mathbf{-1} & A_{I_+}^{J_+}
	%	\end{bmatrix}\\&\end{array},$$
	Multiplying by $\diag(-A^{j+1})$ and $\diag(A_{i+1})$, we obtain the following:
	$$\diag(-A^{j+1})A\diag(A_{i+1}) = \begin{array}{c@{}c}
& \begin{array}{cccc}
	1 & J_-& j+1 & J_+
\end{array}\\
\begin{array}{c}
	1\\I_- \\ i+1 \\ I_+
\end{array} &\begin{bmatrix}
	-1 & \mathbf{1}^\top & 1 & \mathbf{-1}^\top \\
	\mathbf{1} & -A_{I_-}^{J_-} & \mathbf{1} & A_{I_-}^{J_+}\\
	1& \mathbf{1}^\top &1&\mathbf{1}^\top\\ 
	\mathbf{-1}& A_{I_+}^{J_-} & \mathbf{1} & -A_{I_+}^{J_+}
\end{bmatrix}\\&\end{array},$$
therefore 
$$\nega (\diag(-A^{j+1})A\diag(A_{i+1})) =\begin{bmatrix}
1 &\mathbf{0}^\top & 0& \mathbf{1}^\top\\
\mathbf{0} & \overline{B_{I_1}^{J_1}} &\mathbf{0}& {B_{I_1}^{J_0}} \\
0 & \mathbf{0}^\top & 0 &\mathbf{0}\\
\mathbf{1} & {B_{I_0}^{J_1}} &\mathbf{0}& \overline{B_{I_0}^{J_0}} \\
\end{bmatrix},$$
with $J_1 = \{j-1\colon j\in J_-\}$, $J_0 = \{j-1\colon j\in J_+\}$, $I_1 = \{i-1\colon i\in I_-\}$, $I_0 = \{i-1\colon i\in I_+\}$.
This yields, $\diag(-A^{j+1})A\diag(A_{i+1})\trim(i+1,j+1) \simeq -2 B_{[i]}^{[j]}$, up to cyclic permutation of the first $i$ rows and $j$ columns.
Finally, by~\eqref{equation:diag_trim_col} \begin{align*}
A\trim(i+1,j+1) &\simeq 2 \diag(\mathbf{1} -2 B^j) B_{[i]}^{[j]}\diag(\mathbf{1}-2 B_i)\\
& \simeq -2 A_{i+1}^{j+1} \diag(\mathbf{1} -2 B^j) B_{[i]}^{[j]}\diag(\mathbf{1}-2 B_i).
\end{align*}


%				Then, if we say $I_1=\supp(B_i)$ and $I_0 =\{1,\ldots,m\}\setminus I_1$, we have up to permutation of the columns 
%				$$B_{[i]} =\renewcommand{\arraystretch}{1.5} \begin{bmatrix}
%					\overline{B_{<i}^{I_1}} & B_{<i}^{I_0} \\
%					\mathbf{1}^\top & \mathbf{0} \\
%					\overline{B_{>i}^{I_1}} & B_{>i}^{I_0} \\					
%				\end{bmatrix},$$




%				$B = \begin{bmatrix}
%					B_1^I & B_1^{I'}\\
%					\vdots & \vdots \\
%					B_{i-1}^I & B_{i-1}^{I'}\\
%					\mathbf{1}^\top & \mathbf{0}^\top \\
%					B_{i+1}^I & B_{i+1}^{I'}\\
%					\vdots & \vdots \\
%					B_{m}^I & B_{m}^{I'}\\
%				\end{bmatrix},\text{ so } B_{[i]} = \begin{bmatrix}
%				\overline{B_1^I} & B_1^{I'}\\
%				\vdots & \vdots \\
%				\overline{B_{i-1}^I} & B_{i-1}^{I'}\\
%				\mathbf{1}^\top & \mathbf{0}^\top \\
%				\overline{B_{i-1}^I} & B_{i-1}^{I'}\\
%				\vdots & \vdots \\
%				\overline{B_{m}^I} & B_{m}^{I'}\\
%				\end{bmatrix}$






\end{proof}

\end{document}
